\documentclass[12pt]{amsart} 
\usepackage{amsfonts}
\usepackage{amscd}
\usepackage{amsmath}
\usepackage{amssymb}
\usepackage{amsthm}
\usepackage{enumitem}
\usepackage{pgf,tikz,tikz-cd}
\usepackage{mathrsfs}
\usepackage[margin=1in]{geometry}
\usepackage{array}
\usepackage{lipsum}
\usepackage{dsfont}
\usepackage{stmaryrd}

%\graphicspath{ {./diagrams/} }

\input{./latex-preamble/cm-preamble.tex}

\newtheorem*{lemma}{Lemma}

\begin{document}

{Atiyah and Macdonald} \hfill {05-19-2026}

\bigskip\bigskip

\begin{center}

{\sc Chapter 3}

\end{center}

\begin{center}

  {Noah Parker}
    
\end{center}

\bigskip

\begin{enumerate}
	\item[16.] {
      Let $B$ be a flat $A$-algebra.
      Then the following conditions are equivalent:
      \begin{enumerate}[label=(\roman*)]
        \item $\a^{ec}=\a$ for all ideals $\a$ of $A$.
        \item $\Spec B\to \Spec A$ is surjective.
        \item For every maximal ideal $\m$ of $A$, we have $\m^e\neq (1)$.
        \item If $M$ is any non-zero $A$-module, then $M_B\neq 0$.
        \item For every $A$-module $M$, the natural mapping $x\mapsto 1\otimes x$ from $M$ to $M_B$ is injective.
      \end{enumerate}
    }
\end{enumerate}

\begin{proof}
  (i)$\iff$(ii): 
  This follows from (3.16).

  (i)$\implies$(iii):
  Suppose not.
  Then we can find a maximal ideal $\m\lideal A$ such that 
  \[
    \m=\m^{ec}=(1)^c=(1),
  \]
  contradicting maximality of $\m$.

  % (iii)$\impliedby$(iv):
  % Suppose that (iv) holds.
  % Let $\m$ be a maximal ideal of $A$, and consider the $A$-module $A/\m\neq 0$.
  % Then (iv) gives us
  % \begin{align*}
  %   0 \neq B\otimes_A (A/\m) \cong \frac{B}{\m^e},
  % \end{align*}
  % so that $\m^e\neq (1)$ as desired.

  (iii)$\implies$(iv):
  Suppose that (iii) holds, and take an $A$-module $M$ such that $M_B=0$.
  We wish to show that $M=0$.
  $B$ is flat, so for any submodule $N\inc M$, we have
  \[
    N_B\cong \im (N_B\inc M_B)\subset M_B =0
  \]
  is zero.
  $M=0$ if any finitely generated submodule $N\inc M$ is zero, so we may assume that $M$ is finitely generated.

  Suppose, by way of contradiction, that $M\neq 0$.
  (3.8) gives us $(M_B)_\m= 0$ for all maximal ideals $\m$ of $A$, as well as some maximal ideal $\m$ with $M_\m\neq 0$.
  Fix such a maximal ideal $\m$ with $M_\m\neq 0$, and let $K=A_\m/\m A_\m$.
  Then
  \begin{align*}
    0 &= (M_B)_K \\
      &= K\otimes_A B\otimes_A M \\
      &= B_K\otimes_K M_K
  \end{align*}
  implies $B_K=0$ or $M_K=0$, since $K$ is a field.
  Nakayama's lemma gives us $M_K\neq 0$, so we must have $B_K=0$.

  Let $k=A/\m$, so $\m=(\m A_\m)^c$ implies $k\inc K$ is a field extension.
  $B$ flat gives us that
  \begin{align*}
    \frac{B}{\m^e B}\cong B_k\inc B_K=0
  \end{align*}
  is an inclusion.
  Then $B/\m^e B$ implies $\m^e=(1)$, contradicting our assumption.
  We conclude that, in fact, $M=0$.

  % (iv)$\impliedby$(v): 
  % This is clear.

  (iv)$\implies$(v):
  Suppose (iv) holds, let $M$ be an $A$-module, and let $x\in \ker(M\to M_B)$.
  $B$ flat gives us that $(x)_B\inc M_B$ is injective, so that $(x)_B\cong (1\otimes x)=0$ implies $(x)=0$ by (iv).
  Then $x=0$ gives us $\ker(M\to M_B)=0$, so that (v) holds.

  (v)$\implies$(i):
  Suppose (v) holds, and let $\a$ be an ideal of $A$.
  Using (v) and flatness of $B$, we can write down the following diagram.
  \[
    \begin{tikzcd}
      A\ar[r,hook] & B \\
      \a\ar[u,hook]\ar[r,hook]& \a_B\ar[u,hook]
    \end{tikzcd}
  \] 
  This allows us to identify 
  \begin{align*}
    \a=\im(\a\inc A)=\im(\a_B\inc B)^c=\a^{ec},
  \end{align*}
  as desired.
\end{proof}

\begin{enumerate}
  \item[17.] {
      Let $A\xrightarrow{f}B\xrightarrow{g}C$ be ring homomorphisms.
      If $g\circ f$ is flat and $g$ is faithfuly flat, then $f$ is flat.
    }
\end{enumerate}

\begin{proof}
  Let $N\inj M$ be a monomorphism of $A$-modules.
  The following diagram commutes.
  \[
    \begin{tikzcd}
      N_C\ar[r,tail]&M_C \\
      N_B\ar[u,tail]\ar[r]&M_B\ar[u,tail] \\
      N\ar[u]\ar[r,tail]&M\ar[u]
    \end{tikzcd}
  \] 
  Then
  \begin{align*}
    \ker(N_B\to M_B)&\subset \ker(N_B\to M_B\to M_C) \\
                    &=\ker(N_B\to n_C\to M_C) 
                    =0
  \end{align*}
  gives us injectivity of $N_B\to M_B$, so that $f$ is indeed flat.
\end{proof}

\end{document}
